% Volume VIII: Applications and Case Structures
% Part XV. Domains of Difficult Judgment
% Chapter 84: Public Disorder and Mass Recording
% Spine: proof-spines/ch084-spine.md

\chapter{Public Disorder and Mass Recording}
\label{ch:ch084}

Events involving thousands of participants and millions of recordings provide the clearest example of evidence abundance without rapid comprehension. This chapter studies public disorder through multi-camera reconstruction, identity uncertainty, decentralized action, provocation, and temporal fragmentation. The central problem is that recordings proliferate faster than they can be synchronized, authenticated, and interpreted, so visibility expands more quickly than understanding.

The task, then, is not to deny what cameras reveal but to specify what each recording can bear. A clip may establish presence while leaving role, sequence, intent, and instigation unresolved; a large archive may promise total legibility while remaining practically unmastered in real time. The chapter uses this setting to show why contested crowd events require methods of reconstruction that match the scale and fragmentation of the evidence itself.
